%=============================================================================
\section{A Topological Link Between $H_0$ and $M_P$}\label{sec:G-derivation}
%=============================================================================

The preceding sections treated $M_P$ (equivalently $G$) as an input parameter. Here we present a dimensionless constraint linking $G$, $\hbar$, $H_0$, $c$, and $\alpha$, such that \textit{given one scale measurement, all others follow from topology}.

%-----------------------------------------------------------------------------
\subsection{The Dimensionless Invariant}
%-----------------------------------------------------------------------------

The primary claim is a purely dimensionless relation, now closed to theorem status via the observer dictionary (Appendix~\ref{app:O_alpha57_clean}):

\begin{proposition}[Topological Invariant --- Dictionary-Closed]\label{prop:invariant}
DFD predicts the following dimensionless constraint:
\begin{equation}
\boxed{\frac{G\hbar H_0^2}{c^5} = \alpha^{k_{\max} - N_{\rm gen}} = \alpha^{57}}
\label{eq:invariant}
\end{equation}
where $k_{\max} = 60$ (Spin$^c$ index from Lemma~\ref{lem:E_bundle_kmax60}), $N_{\rm gen} = 3$ (generation count), and $\alpha$ is the fine-structure constant.

\textbf{Theorem-grade status (Appendix~\ref{app:O_alpha57_clean}):}
\begin{itemize}
\item The exponent $57 = k_{\max} - N_{\rm gen}$ is forced by primed-determinant scaling on the finite Toeplitz state space (Lemma~\ref{lem:O_det_scaling}, Corollary~\ref{cor:O_alpha57}).
\item The identification with the observed invariant $I = G\hbar H_0^2/c^5$ is made explicit as the observer dictionary (Definition~\ref{def:O_dictionary}).
\end{itemize}
\end{proposition}

This formulation has several advantages:
\begin{itemize}
\item \textbf{Dimensionless}: No unit conventions or hidden factors
\item \textbf{Symmetric}: Predicts $G$ from $H_0$ \textit{or} $H_0$ from $G$
\item \textbf{Falsifiable}: A single testable constraint
\end{itemize}

\paragraph{Bidirectionality.}
Given $(\alpha, \hbar, c)$ and a measured $G$, the invariant predicts $H_0$. Equivalently, given $H_0$ it predicts $G$. Neither is privileged as ``input''---the constraint is symmetric. This prevents any accusation that one quantity was ``chosen'' to match the other.

\paragraph{Error propagation.}
Taking logarithms and differentiating:
\begin{equation}
\frac{\delta G}{G} = -2\frac{\delta H_0}{H_0}
\label{eq:error-prop}
\end{equation}
The precision of any $G$ prediction is limited by $H_0$ uncertainty. With current $H_0$ uncertainties of $\sim$1--2\%, the constraint tests $G$ at the $\sim$2--4\% level.

\paragraph{Equivalent form (Planck mass).}
Defining $M_P = \sqrt{\hbar c/G}$, the invariant becomes:
\begin{equation}
M_P = \alpha^{-(k_{\max} - N_{\rm gen})/2} \times \frac{\hbar H_0}{c^2} = \alpha^{-28.5} \times \frac{\hbar H_0}{c^2}
\label{eq:MP-master}
\end{equation}

\paragraph{Numerical verification.}
Using CODATA values for $G$, $\hbar$, $c$, $\alpha$:
\begin{align}
\text{LHS: } \frac{G\hbar H_0^2}{c^5} &= 1.587 \times 10^{-122} \quad (\text{at } H_0 = 72.1~\text{km/s/Mpc}) \nonumber\\
\text{RHS: } \alpha^{57} &= 1.586 \times 10^{-122}
\end{align}
Agreement to 0.03\% on a quantity spanning 122 orders of magnitude.

%-----------------------------------------------------------------------------
\subsection{Implication for the Cosmological Constant Problem}
%-----------------------------------------------------------------------------

The cosmological constant problem asks: why is $\rho_\Lambda / \rho_{\rm Planck} \approx 10^{-123}$? This is often called ``the worst fine-tuning in physics'' because naive quantum field theory predicts $\rho_\Lambda \sim \rho_{\rm Planck}$.

If Eq.~\eqref{eq:invariant} holds, the ratio is \textit{topologically constrained}:

\begin{proposition}[Cosmological Constant Scaling]\label{prop:CC-scaling}
The critical density satisfies:
\begin{equation}
\frac{\rho_c}{\rho_{\rm Planck}} = \frac{3}{8\pi} \times \frac{G\hbar H_0^2}{c^5} = \frac{3}{8\pi}\alpha^{57} \approx 1.9 \times 10^{-123}
\label{eq:CC-solution}
\end{equation}
With $\Omega_\Lambda \approx 0.7$: $\rho_\Lambda/\rho_{\rm Planck} \approx 1.3 \times 10^{-123}$.
\end{proposition}

\begin{proof}[Derivation]
The critical density is $\rho_c = 3H_0^2/(8\pi G)$. The Planck density is $\rho_{\rm Planck} = c^5/(\hbar G^2)$. Thus:
\begin{equation}
\frac{\rho_c}{\rho_{\rm Planck}} = \frac{3H_0^2}{8\pi G} \times \frac{\hbar G^2}{c^5} = \frac{3}{8\pi} \times \frac{G\hbar H_0^2}{c^5}
\end{equation}
Substituting Eq.~\eqref{eq:invariant} gives the result.
\end{proof}

The exponent $57 = k_{\max} - N_{\rm gen} = 60 - 3$ traces to topology:
\begin{itemize}
\item $k_{\max} = 60$: the Spin$^c$ index $\chi(\mathbb{C}P^2, E)$ for twist bundle $E = \mathcal{O}(9) \oplus \mathcal{O}^{\oplus 5}$ (Lemma~\ref{lem:E_bundle_kmax60})
\item $N_{\rm gen} = 3$: the generation count from flux quantization on $S^3$
\end{itemize}

\begin{tcolorbox}[colback=blue!5!white, colframe=blue!60!black, title=Cosmological Constant: Dictionary-Closed Resolution (Appendix~\ref{app:O_alpha57_clean})]
The ``fine-tuning'' of $10^{-123}$ is now closed via the observer dictionary:
\begin{equation}
\frac{\rho_c}{\rho_{\rm Planck}} = \frac{3}{8\pi}\alpha^{k_{\max} - N_{\rm gen}} = \frac{3}{8\pi}\alpha^{57} \approx 10^{-123}
\end{equation}
The exponent 57 is topologically forced by primed-determinant scaling (Corollary~\ref{cor:O_alpha57}). The remaining step is the explicit observer-dictionary identification (Definition~\ref{def:O_dictionary}).
\end{tcolorbox}

%-----------------------------------------------------------------------------
\subsection{Testable Consequence: The Hubble Constant}
%-----------------------------------------------------------------------------

Interpreted as an $H_0$ prediction from $(G, \alpha)$, the invariant Eq.~\eqref{eq:invariant} yields:
\begin{equation}
H_0 = \sqrt{\frac{\alpha^{57} c^5}{G\hbar}} = \frac{\alpha^{28.5}}{t_P}
\label{eq:H0-prediction}
\end{equation}
where $t_P = \sqrt{\hbar G/c^5}$ is the Planck time.

Using CODATA values for $G$, $\hbar$, $c$, $\alpha$:
\begin{equation}
\boxed{H_0^{\rm DFD} = 72.09~\text{km/s/Mpc}}
\label{eq:H0-DFD-value}
\end{equation}

This is a \textbf{zero-parameter prediction}---the value follows entirely from the microsector derivation of $\alpha$ and the topological exponent $57 = k_{\max} - N_{\rm gen}$.

\paragraph{Comparison with observations.}
Recent JWST observations provide high-precision tests of this prediction. Two major collaborations have released results:

\begin{table}[h]
\centering
\caption{Hubble constant: DFD prediction vs.\ observations.}\label{tab:H0-comparison}
\begin{ruledtabular}
\scriptsize
\begin{tabular}{lcccc}
\textbf{Source} & $H_0$ & \textbf{Uncert.} & \textbf{$\Delta/\sigma$} & \textbf{Ref.} \\
\hline
\textbf{DFD prediction} & \textbf{72.09} & (theory) & --- & This work \\
\hline
\multicolumn{5}{c}{\textit{Local distance ladder (JWST)}} \\
SH0ES JWST combined & 72.6 & $\pm 2.0$ & $-0.3\sigma$ & \cite{Riess2024JWST} \\
SH0ES JWST Cepheids & 73.4 & $\pm 2.1$ & $-0.6\sigma$ & \cite{Riess2024JWST} \\
SH0ES JWST TRGB & 72.1 & $\pm 2.2$ & $0.0\sigma$ & \cite{Riess2024JWST} \\
SH0ES JWST JAGB & 72.2 & $\pm 2.2$ & $-0.05\sigma$ & \cite{Riess2024JWST} \\
CCHP TRGB (HST+JWST) & 70.4 & $\pm 1.9$ & $+0.9\sigma$ & \cite{Freedman2024JWST} \\
CCHP JAGB (JWST) & 67.8 & $\pm 2.7$ & $+1.6\sigma$ & \cite{Freedman2024JWST} \\
\hline
\multicolumn{5}{c}{\textit{CMB-inferred (model-dependent)}} \\
Planck $\Lambda$CDM & 67.4 & $\pm 0.5$ & $+9.4\sigma$ & \cite{Planck2018VI} \\
\end{tabular}
\end{ruledtabular}
\vspace{1mm}
{\scriptsize Units: km/s/Mpc. $\Delta/\sigma \equiv (H_0^{\rm DFD} - H_0^{\rm obs})/\sigma_{\rm obs}$.}
\end{table}

\paragraph{Assessment.}
The DFD prediction $H_0 = 72.09$ km/s/Mpc lies near recent JWST distance-ladder estimates ($\sim$72--73 km/s/Mpc from SH0ES) but above some TRGB/JAGB-based determinations ($\sim$68--70 km/s/Mpc from CCHP). The two JWST teams obtain systematically different results, with the disagreement not yet resolved~\cite{Riess2024JWST,Freedman2024JWST}.

Key observations:
\begin{itemize}
\item The DFD prediction is consistent with all SH0ES JWST measurements within $1\sigma$
\item CCHP results lie $1$--$2\sigma$ below the DFD prediction
\item The Planck CMB-inferred value disagrees at $9.4\sigma$
\end{itemize}

\paragraph{The Hubble tension in DFD.}
The ``Hubble tension''---the $\sim$5 km/s/Mpc discrepancy between local and CMB-inferred values---has a natural interpretation in DFD:
\begin{itemize}
\item \textbf{Local measurements} (Cepheids, SNe Ia) measure actual photon propagation through the $\psi$-field, yielding $H_0 \approx 72$--73 km/s/Mpc
\item \textbf{CMB inference} uses $\Lambda$CDM to extrapolate from $z \sim 1100$, but this model does not account for the $\psi$-screen optical bias (Section~\ref{sec:psi-tomography})
\end{itemize}
The CMB is observed through an accumulated $\Delta\psi \approx 0.30$ (from $\psi$-tomography), which biases distance inferences in the standard framework. The ``tension'' is not a measurement error but a \textbf{model error} in $\Lambda$CDM.

\begin{tcolorbox}[colback=green!5!white, colframe=green!60!black, title=The G-H$_0$ Link: Sharp Prediction]
\textbf{Prediction:} $H_0 = 72.09$ km/s/Mpc (zero free parameters)\\[2pt]
\textbf{Status:} Consistent with SH0ES JWST ($<1\sigma$); above CCHP TRGB/JAGB ($1$--$2\sigma$); incompatible with Planck $\Lambda$CDM ($9.4\sigma$)\\[2pt]
\textbf{Interpretation:} The Hubble tension reflects the $\psi$-screen optical bias ignored by $\Lambda$CDM\\[2pt]
\textbf{Test:} As JWST completes its full Cepheid sample ($\sim$2025--2026), the prediction becomes testable at sub-percent precision
\end{tcolorbox}

%-----------------------------------------------------------------------------
\subsection{Cosmological Evolution of $G$}
%-----------------------------------------------------------------------------

If the topological constraint Eq.~\eqref{eq:invariant} holds at all times, then as $H(t)$ evolves, so must $G(t)$:
\begin{equation}
G(t) = \frac{\alpha^{57} c^5}{\hbar H(t)^2}
\label{eq:G-evolution}
\end{equation}

As the universe expands and $H$ decreases, $G$ \textit{increases}.

\paragraph{Why haven't we detected varying $G$?}
Measurements of $G$ (lunar laser ranging, binary pulsars) use \textit{atomic} references. In DFD, atomic-frame measurements give:
\begin{equation}
G_{\rm atomic} = G_{\rm photon} \times e^{2\psi_{\rm cosmic}}
\end{equation}
If $\psi_{\rm cosmic}$ evolves to compensate for $H$ evolution, then $G_{\rm atomic}$ remains approximately constant while $G_{\rm photon}$ varies.

This is precisely what the cavity-atom LPI test (Section~\ref{sec:cavity}) can detect: the \textit{difference} between photon-frame and atomic-frame measurements of gravitational coupling.

\paragraph{Connection to early universe.}
At the CMB epoch ($z \sim 1100$), $H(z)/H_0 \sim 33000$. In the photon frame:
\begin{equation}
\frac{G(z=1100)}{G_0} = \left(\frac{H_0}{H(z)}\right)^2 \sim 10^{-9}
\end{equation}
Gravity was vastly \textit{weaker} in the early universe (photon frame). This may affect interpretation of BBN and CMB constraints on $G$.

%-----------------------------------------------------------------------------
\subsection{The Parameter Structure}
%-----------------------------------------------------------------------------

If Eq.~\eqref{eq:invariant} holds, DFD has the following structure:

\begin{table}[h]
\centering
\caption{DFD input/output structure.}\label{tab:parameter-count}
\begin{ruledtabular}
\scriptsize
\begin{tabular}{lll}
\textbf{Category} & \textbf{Quantity} & \textbf{Source} \\
\hline
\multirow{3}{*}{Topological} & $k_{\max} = 60$ & $\chi(\mathbb{C}P^2, E)$ \\
 & $N_{\rm gen} = 3$ & Index theorem \\
 & $\alpha^{-1} = 137$ & CS quant. \\
\hline
Observational & $H_0$ or $G$ & Measured \\
\hline
\multirow{5}{*}{Derived} & $G$ or $H_0$ & Eq.~\eqref{eq:invariant} \\
 & $v = 246$ GeV & $M_P \alpha^8 \sqrt{2\pi}$ \\
 & $\rho_c/\rho_{\rm Pl}$ & Eq.~\eqref{eq:CC-solution} \\
 & All masses & $\alpha$-hierarchy \\
 & All mixings & $\mathbb{C}P^2$ geom. \\
\end{tabular}
\end{ruledtabular}
\end{table}

\paragraph{Parameter counting.}
DFD introduces \textbf{no continuous fit parameters}. The discrete topological sector is uniquely determined by Standard Model structure:
\begin{itemize}
\item Hypercharge integrality fixes $q_1 = 3$ (Lemma~\ref{lem:q1})
\item Minimal integer-charge lift gives $\mathcal{O}(9) = L_Y^{\otimes 3}$
\item Five hypercharged chiral multiplet types fix $n = 5$
\item Within $E = \mathcal{O}(a) \oplus \mathcal{O}^{\oplus n}$, minimal-padding uniquely selects $(a,n) = (9,5)$ with $k_{\max} = 60$
\end{itemize}

One scale measurement ($H_0$ or equivalently $G$) determines all dimensionful quantities via the invariant $G\hbar H_0^2/c^5 = \alpha^{57}$.

\begin{tcolorbox}[colback=blue!5!white, colframe=blue!60!black, title=Zero Continuous Parameters --- Dictionary-Closed (Appendix~\ref{app:O_alpha57_clean})]
DFD introduces no continuous fit parameters. Once the discrete topological sector is fixed by Standard Model structure ($k_{\max} = 60$, $N_{\rm gen} = 3$), the exponent in the dimensionless invariant
\begin{equation}
\frac{G\hbar H_0^2}{c^5} = \alpha^{57}
\end{equation}
is topologically forced by primed-determinant scaling (Corollary~\ref{cor:O_alpha57}). The observer dictionary (Definition~\ref{def:O_dictionary}) identifies this with the measured invariant. One scale measurement ($H_0$ or $G$) then fixes all dimensionful quantities.
\end{tcolorbox}

